% =============================================================================
% Figure 4-1: Top-Level System Block Diagram
% =============================================================================
% Headline figure for the SDD. Shows the v1.0 PMVB platform as three layers:
%   1. External consumers (Workstation, MCP client) on the bench network
%   2. Orchestration head (Pi 5) with the VISA/SCPI plane and MCP plane
%      shown as labeled internal bands
%   3. Module fabric: USB hub + Tier 1, Tier 2, Tier 3 envelopes with
%      example front-panel I/O connecting to DUT blocks
%
% Compile:  pdflatex fig_4_1_top_level.tex
% Convert:  pdftocairo -svg fig_4_1_top_level.pdf
% =============================================================================
\documentclass[border=8pt]{standalone}
%
% =============================================================================
% pmvb-figures.sty -- Poor Man's Validation Bench figure style sheet
% =============================================================================
% Defines the house style for all TikZ figures in PMVB design documents.
%
% Usage in figure source:
%     \documentclass[border=4pt]{standalone}
%     \usepackage{pmvb-figures}
%
% Two color modes:
%     \pmvbDarkMode   -- dark background (default; matches SDD HTML theme)
%     \pmvbLightMode  -- white background (for printed portfolio PDFs)
% =============================================================================
\NeedsTeXFormat{LaTeX2e}
\ProvidesPackage{pmvb-figures}[2026/05/06 PMVB figure styles]
%
% --- Required packages ------------------------------------------------------
\RequirePackage{tikz}
% NOTE: dropped the siunitx option. Our figures use \pmvblabel{} for unit
% annotations rather than circuitikz's SI value formatting, and siunitx.sty
% isn't always available on bare TeX Live installs.
\RequirePackage[american]{circuitikz}
\RequirePackage{xcolor}
\RequirePackage{etoolbox}
%
% --- TikZ libraries ---------------------------------------------------------
\usetikzlibrary{arrows.meta}
\usetikzlibrary{positioning}
\usetikzlibrary{shapes.geometric}
\usetikzlibrary{shapes.misc}
\usetikzlibrary{fit}
\usetikzlibrary{backgrounds}
\usetikzlibrary{calc}
\usetikzlibrary{decorations.pathreplacing}
\usetikzlibrary{patterns}
%
% =============================================================================
% Color palette (FMCW dark theme, matches SDD HTML)
% =============================================================================
\definecolor{pmvbBg}{HTML}{0d1117}
\definecolor{pmvbFg}{HTML}{e6edf3}
\definecolor{pmvbMuted}{HTML}{94a3b8}
\definecolor{pmvbBlue}{HTML}{3b82f6}
\definecolor{pmvbBlueFill}{HTML}{1f2937}
\definecolor{pmvbGreen}{HTML}{10b981}
\definecolor{pmvbGreenFill}{HTML}{1e293b}
\definecolor{pmvbAmber}{HTML}{f59e0b}
\definecolor{pmvbRed}{HTML}{ef4444}
\definecolor{pmvbViolet}{HTML}{a78bfa}
%
\definecolor{pmvbBgLight}{HTML}{ffffff}
\definecolor{pmvbFgLight}{HTML}{0f172a}
\definecolor{pmvbMutedLight}{HTML}{475569}
\definecolor{pmvbBlueFillLight}{HTML}{dbeafe}
\definecolor{pmvbGreenFillLight}{HTML}{d1fae5}
%
% =============================================================================
% Mode selectors
% =============================================================================
\newcommand{\pmvbDarkMode}{%
  \pagecolor{pmvbBg}%
  \color{pmvbFg}%
}
%
\newcommand{\pmvbLightMode}{%
  \colorlet{pmvbBg}{pmvbBgLight}%
  \colorlet{pmvbFg}{pmvbFgLight}%
  \colorlet{pmvbMuted}{pmvbMutedLight}%
  \colorlet{pmvbBlueFill}{pmvbBlueFillLight}%
  \colorlet{pmvbGreenFill}{pmvbGreenFillLight}%
  \pagecolor{pmvbBgLight}%
  \color{pmvbFgLight}%
}
%
\AtBeginDocument{\pmvbDarkMode}
%
% =============================================================================
% Node styles
% -----------------------------------------------------------------------------
% IMPORTANT: each \tikzset{...} block below must be a single argument with NO
% blank lines inside, otherwise \par is inserted and pgfkeys errors out.
% =============================================================================
%
% --- Subsystem (system-level block: a board, a chip, a whole module) ---
\tikzset{
  pmvb subsystem/.style={
    rectangle,
    rounded corners=2pt,
    draw=pmvbBlue,
    line width=0.8pt,
    fill=pmvbBlueFill,
    text=pmvbFg,
    minimum width=28mm,
    minimum height=12mm,
    align=center,
    font=\sffamily\small,
    inner sep=4pt,
  },
}
%
% --- IC / hardware part (functional-level block) ---
\tikzset{
  pmvb ic/.style={
    rectangle,
    rounded corners=1pt,
    draw=pmvbGreen,
    line width=0.6pt,
    fill=pmvbGreenFill,
    text=pmvbFg,
    minimum width=24mm,
    minimum height=10mm,
    align=center,
    font=\sffamily\footnotesize,
    inner sep=3pt,
  },
}
%
% --- Internal block (inside an IC's functional block diagram) ---
\tikzset{
  pmvb internal/.style={
    rectangle,
    draw=pmvbFg,
    line width=0.5pt,
    fill=pmvbBg,
    text=pmvbFg,
    minimum width=18mm,
    minimum height=7mm,
    align=center,
    font=\sffamily\scriptsize,
    inner sep=2pt,
  },
}
%
% --- Pin pad (the box-with-X at IC die boundary in datasheet diagrams) ---
\tikzset{
  pmvb pin/.style={
    rectangle,
    draw=pmvbFg,
    line width=0.4pt,
    fill=pmvbBg,
    minimum size=2.5mm,
    inner sep=0pt,
    path picture={
      \draw[pmvbFg, line width=0.3pt]
        (path picture bounding box.south west) -- (path picture bounding box.north east)
        (path picture bounding box.north west) -- (path picture bounding box.south east);
    },
  },
}
%
% --- Connector (front panel, BNC, banana jack) ---
\tikzset{
  pmvb connector/.style={
    circle,
    draw=pmvbFg,
    line width=0.6pt,
    fill=pmvbBg,
    text=pmvbFg,
    minimum size=8mm,
    font=\sffamily\scriptsize,
    inner sep=1pt,
  },
}
%
% --- Annotation (no border, small muted italic text) ---
\tikzset{
  pmvb annotation/.style={
    text=pmvbMuted,
    font=\sffamily\scriptsize\itshape,
    align=center,
  },
}
%
% --- Subsystem grouping box (used with \node[fit=...]) ---
\tikzset{
  pmvb group/.style={
    rectangle,
    rounded corners=4pt,
    draw=pmvbMuted,
    line width=0.4pt,
    dashed,
    text=pmvbMuted,
    inner sep=8pt,
    font=\sffamily\scriptsize,
  },
}
%
% =============================================================================
% Edge styles
% -----------------------------------------------------------------------------
% Edge styles encode the *kind* of signal flowing on the line.
% =============================================================================
%
\tikzset{
  pmvb digital/.style={
    -Stealth,
    draw=pmvbBlue,
    line width=0.8pt,
    font=\sffamily\scriptsize,
  },
}
%
\tikzset{
  pmvb analog/.style={
    -Stealth,
    draw=pmvbAmber,
    line width=0.8pt,
    font=\sffamily\scriptsize,
  },
}
%
\tikzset{
  pmvb power/.style={
    -Stealth,
    draw=pmvbRed,
    line width=1.0pt,
    font=\sffamily\scriptsize,
  },
}
%
\tikzset{
  pmvb control/.style={
    -Stealth,
    draw=pmvbViolet,
    line width=0.7pt,
    dashed,
    font=\sffamily\scriptsize,
  },
}
%
\tikzset{
  pmvb bus/.style={
    -Stealth,
    draw=pmvbBlue,
    line width=1.4pt,
    double=pmvbBg,
    double distance=0.6pt,
    font=\sffamily\scriptsize,
  },
}
%
% Style for the "xN" bus-width annotation (overlaid on a bus line).
% Dark fill masks the underlying double-line so the white text reads cleanly.
\tikzset{
  pmvb bus width/.style={
    fill=pmvbBg,
    rounded corners=1pt,
    inner sep=2pt,
    font=\sffamily\bfseries\scriptsize,
    text=pmvbFg,
  },
}
%
\tikzset{
  pmvb optional/.style={
    -Stealth,
    draw=pmvbMuted,
    line width=0.6pt,
    dashed,
    font=\sffamily\scriptsize,
  },
}
%
% --- Bidirectional variants ---
\tikzset{
  pmvb digital bi/.style={
    pmvb digital,
    Stealth-Stealth,
  },
}
%
\tikzset{
  pmvb analog bi/.style={
    pmvb analog,
    Stealth-Stealth,
  },
}
%
% =============================================================================
% Diagram-level convenience
% =============================================================================
\tikzset{
  pmvb figure/.style={
    >=Stealth,
    every node/.append style={font=\sffamily},
    node distance=10mm and 14mm,
  },
}
%
% =============================================================================
% Helper macros
% =============================================================================
\newcommand{\pmvblabel}[1]{\textcolor{pmvbMuted}{\sffamily\scriptsize\itshape #1}}
%
\newcommand{\pmvbsubsystem}[3]{%
  node[pmvb subsystem] (#1) {{\bfseries #2}\\\scriptsize #3}%
}
%
\newcommand{\pmvbic}[3]{%
  node[pmvb ic] (#1) {{\bfseries #2}\\\scriptsize #3}%
}
%
% NOTE on bus-width annotations: TikZ's path parser doesn't reliably expand
% \newcommand macros at mid-path positions, so for the "xN" overlay use the
% inline node-in-path syntax directly:
%     \draw[pmvb bus] (A) -- node[pmvb bus width]{$\times 16$} (B);
% The pmvb bus width style is defined above.
%
% -----------------------------------------------------------------------------
% \opamptri{name}{coord}{dir}
%   Draws an op-amp triangle. Apex points right (dir=1) or left (dir=-1).
%   Defines coordinates: <name>_p, <name>_n, <name>_out for wiring.
%   Place any label as a separate \node call after the macro.
%
%   The {coord} argument MUST include its own parentheses, e.g.
%       \opamptri{bufa}{(0.6, -0.5)}{1}
%       \opamptri{tl072}{($(m1e.center) + (4mm, -14mm)$)}{1}
%   This lets callers use either simple coords or calc expressions.
% -----------------------------------------------------------------------------
\newcommand{\opamptri}[3]{%
  \begin{scope}[shift={#2}, xscale=#3]
    \draw[pmvbFg, line width=0.6pt, fill=pmvbBg]
      (0, 0) -- (
%
\begin{document}
\begin{tikzpicture}[pmvb figure]
  %
  % =========================================================================
  % TOP ROW: external clients on the bench network
  % =========================================================================
  \node[pmvb subsystem, minimum width=34mm, minimum height=18mm]
    (workstn) at (-7, 7)
    {\textbf{Workstation}\\\scriptsize pytest authoring\\\scriptsize Grafana browser};
  %
  \node[pmvb subsystem, minimum width=34mm, minimum height=18mm]
    (mcpcln) at (7, 7)
    {\textbf{MCP Client}\\\scriptsize LLM agent\\\scriptsize bench session};
  %
  % =========================================================================
  % MIDDLE: Pi 5 orchestration head with the two control planes shown as
  % labeled internal bands
  % =========================================================================
  \node[pmvb subsystem, minimum width=104mm, minimum height=38mm]
    (pi5) at (0, 2.5) {};
  %
  % Pi 5 title at the top of the box
  \node[font=\sffamily\bfseries\small, text=pmvbFg]
    at ([yshift=-4mm]pi5.north) {Raspberry Pi 5};
  \node[font=\sffamily\scriptsize, text=pmvbFg]
    at ([yshift=-8mm]pi5.north) {Orchestration Head};
  %
  % VISA/SCPI plane band (left half of Pi 5).
  % Band sized 16mm tall to fit a header line + two-line body without clipping.
  \node[pmvb internal, minimum width=46mm, minimum height=16mm,
        anchor=north west]
    at ([xshift=4mm, yshift=-13mm]pi5.north west)
    (visa_plane) {};
  \node[font=\sffamily\bfseries\scriptsize, text=pmvbBlue,
        anchor=north]
    at ([yshift=-1.5mm]visa_plane.north) {VISA/SCPI plane};
  \node[font=\sffamily\scriptsize, text=pmvbFg, align=center,
        anchor=south]
    at ([yshift=1.5mm]visa_plane.south)
    {pytest \textbullet{} PyVISA\\InfluxDB \textbullet{} Grafana};
  %
  % MCP plane band (right half of Pi 5)
  \node[pmvb internal, minimum width=46mm, minimum height=16mm,
        anchor=north east, draw=pmvbViolet]
    at ([xshift=-4mm, yshift=-13mm]pi5.north east)
    (mcp_plane) {};
  \node[font=\sffamily\bfseries\scriptsize, text=pmvbViolet,
        anchor=north]
    at ([yshift=-1.5mm]mcp_plane.north) {MCP plane};
  \node[font=\sffamily\scriptsize, text=pmvbFg, align=center,
        anchor=south]
    at ([yshift=1.5mm]mcp_plane.south)
    {MCP gateway\\(LLM-callable tools)};
  %
  % =========================================================================
  % LAN / wired Ethernet from external clients to Pi 5 planes
  % L-shaped routing: Workstation goes DOWN then RIGHT (entering visa_plane
  % from its left edge); MCP Client goes DOWN then LEFT (entering mcp_plane
  % from its right edge).
  % =========================================================================
  \draw[pmvb digital bi]
    (workstn.south) --
    (workstn.south |- visa_plane.center) --
    (visa_plane.west);
  %
  \draw[pmvb digital bi]
    (mcpcln.south) --
    (mcpcln.south |- mcp_plane.center) --
    (mcp_plane.east);
  %
  % LAN annotation: just above Pi 5's top edge, tight to the box
  \node[font=\sffamily\scriptsize\itshape, text=pmvbMuted]
    at (0, 4.6) {wired Ethernet (bench network)};
  %
  % =========================================================================
  % Pi 5 to USB Hub
  % =========================================================================
  \node[pmvb subsystem, minimum width=70mm, minimum height=14mm]
    (hub) at (0, -2)
    {\textbf{Powered USB Hub}\\\scriptsize 4--8 ports \textbullet{} 5\,V at 2.4\,A/port};
  %
  \draw[pmvb digital bi]
    (pi5.south) --
    node[right=0.5mm]{\pmvblabel{USB 2.0}}
    (hub.north);
  %
  % =========================================================================
  % Tier envelopes (T1, T2, T3 side by side)
  % =========================================================================
  %
  % --- Tier 1 envelope ---
  \node[pmvb subsystem, minimum width=64mm, minimum height=46mm]
    (t1) at (-7.5, -8) {};
  \node[font=\sffamily\bfseries\small, text=pmvbFg]
    at ([yshift=-4mm]t1.north) {Tier 1};
  \node[font=\sffamily\scriptsize, text=pmvbFg]
    at ([yshift=-8mm]t1.north) {MCU-based \textbullet{} $\times 8$};
  \node[pmvb ic, minimum width=46mm, minimum height=10mm]
    at ([yshift=2mm]t1.center) (t1_pico)
    {\textbf{Pico 2 W}\\\scriptsize + analog/digital front end};
  %
  % --- Tier 2 envelope ---
  \node[pmvb subsystem, minimum width=64mm, minimum height=46mm]
    (t2) at (0, -8) {};
  \node[font=\sffamily\bfseries\small, text=pmvbFg]
    at ([yshift=-4mm]t2.north) {Tier 2};
  \node[font=\sffamily\scriptsize, text=pmvbFg]
    at ([yshift=-8mm]t2.north) {FPGA-based \textbullet{} $\times 5$};
  \node[pmvb ic, minimum width=46mm, minimum height=8mm]
    at ([yshift=4mm]t2.center) (t2_pico) {Pico 2 W};
  \node[pmvb ic, minimum width=46mm, minimum height=8mm]
    at ([yshift=-7mm]t2.center) (t2_fpga) {Tang Primer 25K};
  \draw[pmvb digital]
    (t2_pico) -- node[right]{\pmvblabel{SPI 30 MHz}} (t2_fpga);
  %
  % --- Tier 3 envelope (deferred, dashed gray) ---
  \node[pmvb subsystem, minimum width=64mm, minimum height=46mm,
        draw=pmvbMuted, dashed]
    (t3) at (7.5, -8) {};
  \node[font=\sffamily\bfseries\small, text=pmvbMuted]
    at ([yshift=-4mm]t3.north) {Tier 3};
  \node[font=\sffamily\scriptsize\itshape, text=pmvbMuted]
    at ([yshift=-8mm]t3.north) {deferred \textbullet{} Phase 4};
  \node[font=\sffamily\scriptsize, text=pmvbMuted, align=center]
    at (t3.center)
    {Tang Mega 138K Pro\\or AX7325B\\$+$ Pi Zero streaming sidecar};
  %
  % =========================================================================
  % USB-TMC bus: hub -> junction -> three tier envelopes
  % =========================================================================
  \coordinate (bus_jct) at ($(hub.south) + (0, -2.4)$);
  %
  \draw[pmvb bus, -]
    (hub.south) --
    node[pmvb bus width]{$\times 8$}
    (bus_jct);
  %
  \draw[pmvb bus] (bus_jct) -| (t1.north);
  \draw[pmvb bus] (bus_jct) -- (t2.north);
  % Tier 3 deferred: drop the arrowhead so the dashed muted tail reads as a
  % placeholder rather than an active signal connection.
  \draw[pmvb bus, draw=pmvbMuted, dashed, -] (bus_jct) -| (t3.north);
  %
  \node[font=\sffamily\scriptsize\itshape, text=pmvbMuted]
    at ($(hub.south) + (1.6, -1.2)$) {USB-TMC};
  %
  % =========================================================================
  % Front-panel I/O: BNC directly below each tier (centered), DUT directly
  % below the BNC. Internal wire from the front-end IC down through the
  % tier south edge to the BNC, then to the DUT.
  % =========================================================================
  %
  % --- Tier 1 BNC + DUT ---
  \node[pmvb connector, scale=0.7]
    at ($(t1.south) + (0, -7mm)$) (t1_bnc) {BNC};
  \node[pmvb subsystem, minimum width=24mm, minimum height=12mm,
        fill=pmvbBg, draw=pmvbMuted]
    at ($(t1_bnc.south) + (0, -1.0)$) (dut1) {DUT};
  %
  \draw[pmvb analog] (t1_pico.south) -- (t1_bnc.north);
  \draw[pmvb analog] (t1_bnc.south) -- (dut1.north);
  %
  % --- Tier 2 BNC + DUT ---
  \node[pmvb connector, scale=0.7]
    at ($(t2.south) + (0, -7mm)$) (t2_bnc) {BNC};
  \node[pmvb subsystem, minimum width=24mm, minimum height=12mm,
        fill=pmvbBg, draw=pmvbMuted]
    at ($(t2_bnc.south) + (0, -1.0)$) (dut2) {DUT};
  %
  \draw[pmvb analog] (t2_fpga.south) -- (t2_bnc.north);
  \draw[pmvb analog] (t2_bnc.south) -- (dut2.north);
  %
  % =========================================================================
  % Optional trigger bus: VERTICAL line in the gap between T1 and T2,
  % with horizontal taps going LEFT into T1.east and RIGHT into T2.west
  % at the tier center y. Annotation sits below, in the open area between
  % DUT1 and DUT2.
  % =========================================================================
  \coordinate (trig_x)      at ($(t1.east)!.5!(t2.west)$);
  \coordinate (trig_top)    at (trig_x |- t1.north);
  \coordinate (trig_bottom) at (trig_x |- t1.south);
  \coordinate (trig_tap)    at (trig_x |- t1.center);
  %
  \draw[pmvb optional, -] (trig_top) -- (trig_bottom);
  \draw[pmvb optional] (trig_tap) -- (t1.east);
  \draw[pmvb optional] (trig_tap) -- (t2.west);
  %
  \node[font=\sffamily\scriptsize\itshape, text=pmvbMuted, align=center]
    at ($(trig_bottom) + (0, -8mm)$)
    {Optional shared\\trigger bus};
  %
  % =================================